\documentclass[exercice,a4]{thedocument}%
\usepackage{texmaths-tikz}%
\usepackage{pstricks}%
\usepackage{pst-bezier}%
\usepackage{pstricks-add}%
\usepackage{pst-eucl}%
\usepackage{pst-func}%
\usepackage{pst-tree}%
\usepackage{pst-circ}%
\usepackage{tikz}%
\usepackage{pgfplots}%
\usepackage{tkz-tab}%
\usepackage{tikzpagenodes}%
\usepackage{tabularray}%
\usepackage{tabularx}%
\usepackage{multirow}%
\usepackage{multicol}%
\usepackage{scratch3}%
\usepackage{enumitem}%
\usepackage[np]{numprint}%
\usepackage{pas-tableur}%
\usepackage{listings}%
\usepackage{graphicx}%
\usepackage{adjustbox}%
\usepackage{needspace}%
\usepackage{xcolor}%
\usepackage{eurosym}%
\usepackage{diagbox}%
\usepackage{colortbl}%
\usepackage{multido}%
\usepackage{icomma}%
\usepackage[french]{babel}%
\startdatakeys
    \source{26GENMATPO3}
    \BO{2026}
    \cycle{4}
    \competencesDuSocle{CA,RA,CO}
    \objectifsApprentissage{C4A1aa,C4A2ad,C4B2cb,C4D2aa,C4D3ac}
    \motsCles{symétrie axiale,Pythagore,Thalès,construction,droites parallèles}
\enddatakeys
\definecolor{gris}{gray}{0.5}%
\providecommand{\wideparen}[1]{\overset{\frown}{#1}}%
\begin{document}%
\providecommand{\degres}{^\circ}%
On considère la figure géométrique ci-dessous qui n'est pas en vraie grandeur.

Les points G, A, O, C et E sont alignés.

Les points B, O, D et F sont alignés.

COD est un triangle rectangle en O.

$\mathrm{OD} = 4$ cm, $\mathrm{OB} = 6$ cm, $\mathrm{OC} = 3$ cm, $\mathrm{CE} = 1{,}5$ cm et $\mathrm{DF} = 2$ cm.

La droite (BF) est l'axe de symétrie de cette figure.

\begin{center}
\begin{tikzpicture}[rotate=25,scale=1.2, line cap=round, line join=round, font=\footnotesize]
	\coordinate (O) at (0,0);
	\coordinate (C) at (2,0);
	\coordinate (A) at (-2,0);
	\coordinate (D) at (0,-2.5);
	\coordinate (B) at (0,3.5);
	\coordinate (E) at (3,0);
	\coordinate (F) at (0,-3.8);
	\coordinate (G) at (-3,0);
	\draw[line width=1.2pt] (A) -- (B) -- (C) -- (D) -- cycle;
	\draw[line width=1.2pt] (A) -- (G) -- (F) -- (D);
	\draw[line width=1.2pt] (C) -- (E) -- (F);
	\draw[line width=0.8pt] (G) -- (E);
	\draw[line width=0.8pt] (B) -- (F);
	\node[above left] at (A) {A};
	\node[above] at (B) {B};
	\node[below right] at (C) {C};
	\node[below right] at (D) {D};
	\node[right] at (E) {E};
	\node[below] at (F) {F};
	\node[left] at (G) {G};
	\node[below left] at (O) {O};
	\draw[red, line width=1pt] (0,0) rectangle (0.3,-0.3);
	
	\node at (0.8,0.3) {$3$ cm};
	\node at (2.4,0.3) {$1,5$ cm};
	\node at (-0.3,1.6) {$6$ cm};
	\node at (-0.3,-1.1) {$4$ cm};
	\node at (-0.3,-3) {$2$ cm};
\end{tikzpicture}
\end{center}

\begin{enumerate}
	\item Justifier que $\mathrm{OA} = 3$ cm.
	\item Sur la copie, construire en vraie grandeur le quadrilatère ABCD.
	\item En se plaçant dans le triangle OFE, calculer la longueur du côté [FE].
	\item Démontrer que les droites (CD) et (EF) sont parallèles.
\end{enumerate}

\needspace{10\baselineskip}
\medskip
\showthesource
\end{document}
